\section{Two inherited algorithmic obligations reconstructed}\label{sec:algorithmic}
The following local results have complete written proofs. They do not
settle the separate alias-estimator or finite-precision arguments.

\begin{theorem}[Transformation and isolated bin]\label{thm:isolation}
Let $a_i$ be positive integers and $t\ge0$. Choose a power of two
$N_0\in[12n,24n)$, independent uniform $u_i\in\{1,\ldots,N_0\}$,
and put $C=1+nN_0$, $c_i=Ca_i+u_i$, $t_s=Ct+s$ for
$0\le s\le nN_0$. For every choice of the tie weights,
\[
 \exists I\subseteq[n]:a(I)=t
 \quad\Longleftrightarrow\quad
 \exists s,I:c(I)=t_s.
\]
On a NO input every transformed bin is empty. On a YES input the
probability that at least one transformed bin has exactly one index-set
representation is at least $11/12$.
\end{theorem}
\begin{proof}
If $a(I)=t$, take $s=u(I)\in[0,nN_0]$. Conversely,
$c(I)=Ct+s$ implies $C(a(I)-t)=s-u(I)$, whose absolute value is at
most $nN_0<C$. Since $a(I)-t$ is an integer, it is zero.

For completeness, let $\mathcal F$ be any nonempty family of subsets
of $[n]$, here the original solution family. For each coordinate $i$,
fix all $u_j$, $j\ne i$. If both classes exist, let $m_0$ be the
least tie sum among members not containing $i$ and $m_1$ the least
tie sum excluding $u_i$ among members containing $i$. Equality of the
two class minima requires $u_i=m_0-m_1$, at most one of its $N_0$
possible values. If either class is empty there is no such event.
If the overall minimum is attained by two different subsets, some
coordinate distinguishes them, so that coordinate's class minima
are equal. A union bound, conditioning on the other weights for each
coordinate, gives probability at most $n/N_0\le1/12$ of a nonunique
minimum. A unique minimum solution has a unique tie sum within the
solution family, and hence a unique transformed bin by the equivalence
already proved. Nothing excludes repeated $a_i$ or the empty subset.
\end{proof}

\begin{lemma}[Deficiency on uniform half-blocks]\label{lem:deficiency}
For fixed transformed weights define
\[
 \Sigma_c(V)=\{c(I):I\subseteq V\},\qquad
 d_c(V)=|V|-\log_2|\Sigma_c(V)|.
\]
Inserting one index changes deficiency by a value in $[0,1]$.
Replacing one index by another changes it by at most one in absolute
value. For a uniformly sampled $k$-subset $V_k$,
\[
 \operatorname{Var}(d_c(V_k))\le k/4\le n/4.
\]
Also $f(k)=\mathbb E d_c(V_k)$ is nondecreasing in $k$.
\end{lemma}
\begin{proof}
For $i\notin V$,
$\Sigma_c(V\cup\{i\})=\Sigma_c(V)\cup(c_i+\Sigma_c(V))$.
Its cardinality lies between $|\Sigma_c(V)|$ and
$2|\Sigma_c(V)|$, proving the insertion claim. Two sets differing
by an exchange both arise by inserting an index into their common
part, so their deficiencies lie in an interval of length one.
Sampling a uniform $(k+1)$-subset and deleting a uniform member
couples the uniform sizes $k+1,k$; monotonicity follows.

Expose the first $k$ members $X_1,\ldots,X_k$ of a uniformly random
ordered sample without replacement, and let $M_j$ be the conditional
expectation of the final deficiency after $j$ exposures. At a fixed
prefix, compare next choices $a,b$. Transposing $a,b$ gives a bijection
between the possible remaining completions. The corresponding final
sets are either identical (both indices eventually occur) or differ
by one exchange. Hence the two conditional expectations differ by at
most one. Conditional on the prefix, $M_j$ therefore lies in an
interval of length at most one. A real random variable in an interval
of length one has variance at most $1/4$: expand
$(X-\min X)(\max X-X)\ge0$ and complete the square in its mean.
Thus each martingale difference has conditional second moment at most
$1/4$. Distinct differences are orthogonal by iterated conditional
expectation. Summing their variances gives the bound $k/4$.
\end{proof}

\begin{theorem}[A capped preliminary search and the LD split]\label{thm:preliminary}
Let $n\ge2$, $b\ge1$, and use the fixed published $0<\tau<1/2$.
Fix the tie weights and $\zeta>0$. The condition
\begin{equation}\label{eq:LD}
 \mathbb E d_c(V_{\lfloor n/2\rfloor})
       \le(\tfrac12-\tau+\zeta)n
       \tag{LD}
\end{equation}
is exactly the low-deficiency condition used here. There is a
preliminary procedure returning an exact decision or CONTINUE, using
$O(2^{\tau n}\operatorname{poly}(n+b+2))$ word operations on every
history. If LD fails, for all sufficiently large $n$ it returns
CONTINUE with probability at most
\[
 \frac{2}{\zeta^2 n}+n2^{-n^2}.
\]
This conclusion concerns a fixed-input probability; an observed
CONTINUE is not a certificate that LD holds.
\end{theorem}
\begin{proof}
Sample a uniform half-block $V$ and its complement, of sizes
$\lfloor n/2\rfloor$ and $\lceil n/2\rceil$. If LD fails, either
block $U$ has expected deficiency greater than
$(\tfrac12-\tau+\zeta)n$, by Lemma~\ref{lem:deficiency}.
Chebyshev's inequality with variance at most $n/4$ gives
\[
 \Pr\{d_c(U)<(\tfrac12-\tau+\zeta/2)n\}
 \le \frac1{\zeta^2n}.
\]
Outside these two bad events,
\[
 \log_2|\Sigma_c(U)|=|U|-d_c(U)
 \le\tau n-\zeta n/2+1/2<\tau n
\]
when $n>1/\zeta$. Both lists of distinct half sums fit within
$B_n=2^{\lceil\tau n\rceil}$.

Construct such a sorted list by starting with $\{0\}$ and, for each
weight, merging the current sorted list with its translate, deleting
duplicate values. If the next distinct output would exceed $B_n$,
return CONTINUE before writing it. The number of distinct sums is
nondecreasing under insertions, so if the final list fits, every prefix
fits. Each merge costs $O(B_n)$ operations. Once both full lists fit,
the sorted two-pointer search answers exactly whether some left and
right sums total any required $t_s$. There are $O(n^2)$ targets, so
even checking them separately costs only a polynomial times $B_n$.
If either list hits its cap, no inexact decision is emitted.

For an explicit bounded sampler, expose a uniform ordered sample
without replacement using partial Fisher--Yates swaps. A required
integer in $\{0,\ldots,m-1\}$ is obtained by masking a fresh uniform
word to $\lceil\log_2m\rceil$ bits and rejecting values at least $m$.
Acceptance is at least $1/2$. Cap each draw at $n^2$ proposals and
return CONTINUE on exhaustion. Couple this with the ideal unbounded
sampler: at most $n$ draws are made, so discrepancy probability is at
most $n2^{-n^2}$. The ideal subset is exactly uniform. The bounded
sampler has polynomial cost on every history.

All sums have $O(b+\log n)$ bits: $N_0<24n$ and $C<24n^2+1$
give $\sum_i c_i<31n^3 2^b$ for $n\ge2$, $b\ge1$.
The bit bound $b+3\log_2n+5$ is below the prescribed
$\lceil4(n+b+\log_2(n+2))\rceil$ word size on this domain, so sums
and keys occupy a fixed number of words. Two capped final lists plus
the alternating merge workspace can be kept in at most $4(B_n+1)$
words: the translated list is read implicitly, not copied. The partial
Fisher--Yates array, input and counters require at most $K n$ additional
words for a fixed implementation constant $K$. Thus simultaneous
storage is at most $4(B_n+1)+K n$, whose address length is
$\tau n+O(\log n)$. Since $\tau<1/2$, this is below the prescribed
word size for all $n$ above a fixed implementation cutoff. For the
finitely many smaller $n\ge2$, enumerate subsets with a running sum
in $O(n)$ storage and answer exactly; its $2^n$ time contributes only
a fixed factor to the stated $O(2^{\tau n}\operatorname{poly}(n+b+2))$
bound. This cutoff is independent of the input bit length and of $c$.
Fixed rational $\tau$ permits exact calculation of its ceiling.
Every allocation and loop has an explicit cap,
including the failure histories, proving the stated all-history cost.
\end{proof}

These proofs reconstruct the transformation and preliminary-search
results of the pinned OpenAI setup and Kapila Sections~2.1/4
\cite{openai,kapila}. They establish the local statements under their
displayed hypotheses without assuming an exact subset-sum solver or
an LD oracle. In the composed analysis one fixes the tie weights first,
splits according to LD, and draws the main procedure's independent
randomness whether or not an earlier exact answer makes it unused.
This avoids conditioning the main proof on the observed continuation.

\section{The consequential algorithmic dependencies that remain}\label{sec:remaining}
\begin{proposition}[Exact alias classes and sound insertion]\label{prop:alias-sound}
Use the actual integer key $(c(A),b(A))$ for each half-block subset,
where $b_i=c_i\bmod P$ is the least nonnegative residue. Keep a
dictionary key $(\kappa_L,\kappa_R,Z)$ only once, and store the value-map
phase sums $\Phi_L,\Phi_R$ over all subsets with the respective key.
Every insertion made by Kapila's Section~5.2 matching/expansion rule
is an alias class. If all alias classes are reached, the exact output
equals the full alias phase sum, with no multiplicity error from
repeated discoveries or repeated original weights.
\end{proposition}
\begin{proof}
For precision, $T_{s,q}=(t_s\bmod P)+qP$, $0\le q<n$.
Every modular representation has exactly one such carry: its integer
sum of reduced weights is in $[0,n(P-1)]\subset[0,nP)$ and is congruent
to the least residue of $t_s$. Subtracting that residue and dividing
by $P$ gives an integer in $\{0,\ldots,n-1\}$.
Write $c^F_i=(-1)^{[i\in F]}c_i$ on the large block, and similarly
for $b^F$. Directly partitioning $Z$ into $Z\setminus F$ and
$Z\cap F$ proves
$c^F(Z)=c(Z\triangle F)-c(F)$, and the same identity for $b$.
In an insertion the masks $X,Y$ are disjoint; put
$Z=(X\cup Y)\triangle F$. The integer matching condition is
\[
 b(\kappa_L)+b^F(X)+b(\kappa_R)+b^F(Y)
       =T_{s,q}-b(F),
\]
where $T_{s,q}\equiv t_s\pmod P$. The sign identity gives
$b(A\cup Z\cup B)=T_{s,q}$ for every half-block pair having those
keys, and thus its unreduced sum is congruent to $t_s$ modulo $P$.
The skipped exact interval in the matching rule excludes precisely
\[
 c(\kappa_L)+c^F(X)+c(\kappa_R)+c^F(Y)=t_s-c(F).
\]
The same identity therefore shows that each inserted member has
unreduced sum different from $t_s$. It is an alias, not an exact
representation.

Every index set $I$ belongs to exactly one class:
$(\kappa(I\cap L),\kappa(I\cap R),I\cap M)$, for the fixed disjoint
partition $L,M,R$. Every member of a class has the same unreduced
and reduced sums. Additivity of the retained phases gives
\[
 \sum_{I\text{ in class}}\chi(I)
 =\left(\sum_{A:\kappa(A)=\kappa_L}\chi(A)\right)
  \left(\sum_{B:\kappa(B)=\kappa_R}\chi(B)\right)\chi(Z)
 =\Phi_L(\kappa_L)\Phi_R(\kappa_R)\chi(Z).
\]
Distinct class keys partition the relevant index sets, and the
dictionary stores each key at most once. Consequently reaching every
alias class suffices for the claimed exact sum. This does not prove
that the randomized procedure reaches every class before its cap;
that coverage/cost/probability obligation remains below.
\end{proof}

The following are exact obligations, not conclusions of the local
proofs or of scalar certificate checking.
\begin{enumerate}
\item \textbf{Alias coverage and cost.} Under LD, with the published
block sizes and slack, Kapila's Section~5 procedure must return every
modular alias contribution in its deduplicated value classes with the
retained signing, except with a uniformly small failure probability.
Its expected-size estimates must imply the stated deterministic caps
through the actual abort policy, without losing the per-input success
bound. The half-block bad-prime bound uses
$\delta=(\eta-(\mu/2)\log_2 3)/2$; the actual requirements include
$\delta>1/2-\tau+\zeta$ and the bit-length guard slack $\theta<\delta$.
\item \textbf{Same-phase cancellation.} For each target and modular
prime, the alias sum and the checksum must use identical stored phases.
Every modular representation with a different integer sum must occur
exactly once in the deduplicated class sum. This gives exact residual
zero on NO inputs and a modulus-one residual at an isolated bin;
it does not by itself bound either numerical estimator's error.
\item \textbf{Block-table estimator.} For the actual finite-circle
density functions, every selected sampled pattern must satisfy the
table-coverage/second-moment hypotheses in OpenAI Section~4 as adapted
by Kapila Section~6. The released $18\epsilon$ allowance belongs here.
The distribution used to factor coordinate moments must be the original
independent phase law, not a law conditioned on a favorable pattern,
table event, alias-success event or CONTINUE.
\item \textbf{Finite precision and uniform constants.} Scores and
elementary functions must be evaluated with enough accuracy for the
actual block thresholds. Stored phases must remain consistent across
both subroutines; all accumulated errors must stay below the decision
threshold gap. Section~\ref{sec:finite} supplies effective fixed-score
evaluation, but not the full propagated complex-arithmetic error bound.
The program must generate required constants/tables, not receive them
as advice.
\item \textbf{Every-history resources.} Prime search, table construction,
trial lists and visits, rational arithmetic, random draws, memory and
indirect addresses all require deterministic caps with charged word
cost. A finite expected work estimate alone is insufficient. The
published scalar gaps verify proposed exponents, not that a described
procedure actually obeys each one. The guarded program and exact
fallback must be a single program independent of the polynomial-bit
parameter $c$.
\end{enumerate}

\paragraph{Precise cancellation fact.}
If $\chi(I)$ is the same retained unit-modulus phase in both sums,
then, identically,
\[
 \sum_{c(I)\equiv t_s\ (\mathrm{mod}\ P)}\chi(I)
 -\sum_{\substack{c(I)\equiv t_s\ (\mathrm{mod}\ P)\\c(I)\ne t_s}}\chi(I)
 =\sum_{c(I)=t_s}\chi(I).
\]
This is a finite partition identity, valid with repeated original
weights because subsets are index sets. It proves the cancellation
\emph{if} the two enumerated/sampled quantities correspond to those
exact sums with common phases. Completeness of the class algorithm and
accuracy of the estimates remain separate obligations; renaming their
outputs by these sums would assume the desired conclusion.

The current supplement therefore does not claim that the complete
word-RAM algorithm has been independently verified. Any unresolved
point above remains an explicit inherited dependency even after the
boundary, scalar and continuous-moment obligations are settled.

\subsection{A bounded precision obligation that can be discharged locally}
\begin{proposition}[Alias-sum arithmetic, given the correct dictionary]
\label{prop:alias-precision}
Suppose the phase maps and dictionary are those of
Proposition~\ref{prop:alias-sound}, with coverage already supplied.
Use the same exact retained phase vector modulo $P$ throughout.
Compute each unit phase to absolute complex error at most
$\delta_0=2^{-2n^2}$, and round each real and imaginary component
after addition or multiplication to absolute error at most
$R=2^{-n^2}$. Then the total alias-output error is at most
\[
 32\,2^{-n^2+3n}<1/16\qquad(n\ge8).
\]
These operations can be performed with polynomially many charged word
operations per phase or arithmetic operation. This is not a proof of
coverage or of the separate checksum estimator's precision.
\end{proposition}
\begin{proof}
There are at most $2^n$ summands in either half-map and at most $2^n$
dictionary keys: distinct dictionary classes are nonempty and disjoint
subsets of the $2^n$ possible index sets. The Euclidean error of a
complex rounding is at most $2R$. Thus either computed half-map sum
differs from its exact value by at most
$2^n\delta_0+2^{n+1}R\le3\,2^nR=:e$.
The exact sum has modulus at most $2^n$, and its approximation has
modulus at most $2^{n+1}$ for $n\ge2$.

In one class product $\Phi_L\Phi_R\chi(Z)$, replacing the factors
by their approximations changes the value by at most
\[
 e2^n+2^{n+1}e+2^{2n+2}\delta_0
 \le13\,2^{2n}R.
\]
Two rounded multiplications contribute at most $6R$, because the
approximate unit phase has modulus at most two. Hence a class term
has error at most $20\,2^{2n}R$. Summing at most $2^n$ terms and
rounding the additions gives error at most
$20\,2^{3n}R+2^{n+1}R\le32\,2^{3n}R$. The displayed strict
$1/16$ bound holds already at $n=8$ and improves thereafter.
All estimates are deterministic and uniform over the retained phases.

For an explicit effective phase evaluation, first sum the integer
residues modulo $P$ exactly, then approximate the bounded angle
$2\pi k/P$. The identity
$\pi/4=\arctan(1/2)+\arctan(1/3)$ and the alternating arctangent
series give a rational $\pi$ enclosure to error $2^{-8n^2}$ using
$16n^2$ terms. Rational division at the same precision gives an angle
error $O(2^{-8n^2})$. Evaluate sine and cosine by their Taylor
polynomials through degree $16n^2$, on $|x|\le8$. The remainder
is bounded by the exponential-series tail
\[
 \frac{8^{N+1}}{(N+1)!}\,
       \frac1{1-8/(N+2)},\qquad N=16n^2.
\]
For $n\ge8$, use $N!\ge(N/2)^{N/2}$ and $N\ge1024$;
this bound is at most $16\,2^{-24n^2}$.
The angle error and Taylor errors together are strictly less than
$2^{-2n^2}$ in complex norm, with ample slack.

All exact input phases have $O(n)$ bits because the published prime
interval has $\log_2P=O(n)$. Real and imaginary intermediate magnitudes
are at most $2^{O(n)}$, even when summing all dictionary terms.
Keeping $n^2$ fractional bits therefore uses polynomially many bits,
as do the rational Taylor calculations. Elementary multiword integer
arithmetic implements each operation with polynomially many word
operations, charging retained storage and rounding. No transcendental
oracle or uncharged table is required. Finite smaller $n$ can use
the exact fallback already described.
\end{proof}
